% Tvarkaraštis (atsakingas Marijus Planciunas).
% Datos – iš tvarkaraščio modelio (75 veiklos, project-plan/resources/schedule.csv).
% Ganto diagrama (gantt.tex) sugeneruota iš to paties modelio, rankomis jos netaisyti.
% Veiklų ID (A01–A75) komentaruose – tik nuoroda rašančiajam; veiklų lentelės tekste nėra, todėl prozoje
% rašyti veiklos pavadinimą.

\section{Tvarkaraštis}\label{sec:tvarkarastis}

Projektas truks 208~d.~d.: nuo 2026~m. spalio 1~d. iki 2027~m. liepos
30~d. Jį sudaro penki etapai, o Kūrimo etape vyksta keturios iteracijos po
4--5 savaites, kurių kiekviena baigiasi demonstracija užsakovui. Pažangą
žymi 15 tarpinių rezultatų (kontrolinių taškų, KT). Po keturių savaičių
pilotinio naudojimo sistema perduodama užsakovui 2027~m. birželio 29~d.
(KT14), o po to 20~d.~d. trunka garantinis palaikymas.

\subsection{Gyvavimo ciklo modelis ir iteracijos}

Projektas vykdomas pagal UP (angl. \textit{Unified Process}) gyvavimo ciklo
modelį. Jo fazės – pradžia, parengimas, konstravimas ir perdavimas (angl.
\textit{Inception, Elaboration, Construction, Transition}); jų atitiktis
projekto etapams pateikta \ref{tab:etapai}~lentelėje.

Krioklio modelis šiam projektui netinka, nes jame viskas patikrinama tik
pabaigoje. Didžiausi neaiškumai susiję su įranga baseine: ar ji atlaikys
drėgmę, ar radijo ryšys pasieks visus takelius, ar vaizdas ekrane ir garsas
ausinėse sutaps laiku. Pagal krioklio modelį tai paaiškėtų tik baigiamuosiuose
bandymuose, kai visa įranga jau nupirkta ir programa parašyta. Pagal UP
rizikingiausi dalykai tikrinami pirmiausia: bandomasis įrangos komplektas
baseine išbandomas jau Planavimo etape, o grotuvo, sinchronizacijos ir garso
perdavimo prototipai – Projektavimo etape, 2026~m. gruodžio pradžioje. Visa
įranga užsakoma tik po prototipo bandymo baseine.

Kūrimo etapą sudaro keturios iteracijos po 4--5 savaites
(\ref{tab:iteracijos}~lentelė). Kiekviena baigiasi demonstracija užsakovui,
todėl jis kas mėnesį mato veikiančią sistemos dalį. Iteracijos išdėstytos
pagal riziką: pirmosiose kuriama nauja ir techniškai neaiški užsiėmimo
vykdymo grandinė (vykdymo variklis, grotuvas, garsas), o rezervacijos ir
lankomumas – paskutinėje, nes tai įprastas ir gerai žinomas funkcionalumas.
Iteracija prasideda dviejų dienų detaliuoju projektavimu ir baigiasi dviejų
dienų stabilizavimu bei demonstracija; kitos iteracijos projektavimas vyksta
tomis pačiomis dienomis kaip ankstesnės stabilizavimas.

\begin{table}[H]
  \centering
  \caption{Kūrimo etapo iteracijos}
  \label{tab:iteracijos}
  \small
  \begin{tabularx}{\textwidth}{@{}lllr>{\raggedright\arraybackslash}X>{\raggedright\arraybackslash}p{2.1cm}@{}}
    \toprule
    \textbf{Iteracija} & \textbf{Nuo} & \textbf{Iki} & \textbf{D.~d.} & \textbf{Turinys (sistemos dalys)} & \makecell[lb]{\textbf{Tarpinis}\\\textbf{rezultatas}} \\
    \midrule
    1 & 2026-12-04 & 2027-01-11 & 24 & Vykdymo grandinė iš debesies per grotuvą į ekraną, ausines ir kolonėlę (S1, S2, S3, P6, sporto erdvių registras iš P0, S4 prieigos pagrindai) & KT6 \\
    2 & 2027-01-08 & 2027-02-10 & 24 & Užsiėmimų programos su filmuotu turiniu ir darbas be ryšio (P1, P5, S1 atsparumas, S2 talpykla) & KT7 \\
    3 & 2027-02-09 & 2027-03-16 & 24 & Paskyros, užsiėmimų planavimas tvarkaraštyje ir valdymo pultas, istorija ir duomenų apsauga (P0, P2, P7, S4) & KT9 \\
    4 & 2027-03-15 & 2027-04-09 & 19 & Dalyvių valdymas (rezervacijos, lankomumas, priminimai), asmeninės programos ir saugumo stiprinimas (P3, P4, S4) & KT10 (po sistemos bandymų) \\
    \bottomrule
  \end{tabularx}
\end{table}

\subsection{Etapai}

Projektą sudaro penki etapai: Reikalavimų analizė, Planavimas,
Projektavimas, Kūrimas bei Diegimas ir palaikymas (\ref{tab:etapai}~lentelė).
Jų pavadinimai ir datos sutampa su resursų plano laikotarpiais, nes resursų
plane nurodoma, kiek kiekvienos rolės žmonių dirba kiekvienu laikotarpiu;
taip abu skyrius galima tiesiogiai palyginti. Pirmųjų trijų etapų pabaigą
žymi juos užbaigiantys KT: Reikalavimų analizės – KT2, Planavimo – KT3,
Projektavimo – KT4.

Įrangos ir turinio darbai vyksta lygiagrečiai su programinės įrangos kūrimu
ir apima kelis etapus: pratimai filmuojami nuo Reikalavimų analizės etapo,
o įranga sumontuojama Kūrimo etape. Tokia veikla priskiriama etapui, kuriame
ji prasideda. Projekte iš viso dirba 9 žmonės, vienu metu – ne daugiau
kaip 7.

\begin{table}[H]
  \centering
  \caption{Projekto etapai}
  \label{tab:etapai}
  \small
  \begin{tabular}{@{}lllrl@{}}
    \toprule
    \textbf{Etapas} & \textbf{Nuo} & \textbf{Iki} & \textbf{D.~d.} & \textbf{UP fazė} \\
    \midrule
    Reikalavimų analizė & 2026-10-01 & 2026-10-23 & 17 & Pradžia \\
    Planavimas & 2026-10-26 & 2026-11-04 & 7 & Pradžia \\
    Projektavimas & 2026-11-05 & 2026-12-03 & 21 & Parengimas \\
    Kūrimas & 2026-12-04 & 2027-04-09 & 85 & Konstravimas \\
    Diegimas ir palaikymas & 2027-04-12 & 2027-07-30 & 78 & Perdavimas \\
    \midrule
    \textbf{Iš viso} & 2026-10-01 & 2027-07-30 & \textbf{208} & \\
    \bottomrule
  \end{tabular}
\end{table}

\subsection{Tarpiniai rezultatai (kontroliniai taškai)}

Projekto pažangą matuojame ne praleistomis valandomis, o pasiektais
rezultatais. Tam numatyta 15 KT (\ref{tab:kt}~lentelė). Kiekvienas KT turi
kriterijų, kuris nusako, kaip rezultatas patikrinamas: pasirašytas
protokolas, aktas ar užsakovo patvirtinimas. Todėl į klausimą, ar KT
pasiektas, galima atsakyti tik „taip“ arba „ne“. Jei KT vėluoja ar
nepasiektas, projekto vadovas įvertina poveikį ir sprendžia, ar reikia
perplanuoti darbus, ar teikti pakeitimo prašymą.

KT6, KT7, KT9 ir KT10 žymi iteracijų demonstracijas, KT5 ir KT8 –
lygiagrečių turinio ir įrangos darbų pabaigą, KT14 – sistemos perdavimą
užsakovui, o KT15 – projekto uždarymą.

\begin{table}[H]
  \centering
  \caption{Tarpiniai rezultatai (KT)}
  \label{tab:kt}
  \small
  \begin{tabularx}{\textwidth}{@{}l>{\raggedright\arraybackslash}p{0.27\textwidth}l>{\raggedright\arraybackslash}X@{}}
    \toprule
    \textbf{KT} & \textbf{Pavadinimas} & \textbf{Data} & \textbf{Kriterijus} \\
    \midrule
    KT1 & Patvirtintas pratimų ir programų sąrašas & 2026-10-15 & Užsakovo vyriausiasis treneris pasirašė 40~pratimų ir 10~programų sąrašą su aprašais ir filmavimo vieta. \\
    KT2 & Patvirtinta reikalavimų specifikacija & 2026-10-23 & Užsakovas pasirašė funkcinių, nefunkcinių ir aplinkos reikalavimų specifikaciją; atsekamumo lentelėje kiekvienas poreikių aprašo punktas susietas su reikalavimu, atvirų klausimų nėra. \\
    KT3 & Patvirtintas projekto planas ir priėmimo kriterijai & 2026-11-04 & Užsakovas patvirtino tvarkaraštį, iteracijų turinį, rizikų registrą, priėmimo kriterijus ir demonstravimo scenarijus. \\
    KT4 & Patvirtinta architektūra ir baseine patikrintas įrangos sprendimas & 2026-12-03 & Bandomasis komplektas baseine veikė be drėgmės sukeltų gedimų. Ekrane ir dvejose ausinėse vaizdo ir garso sinchronizacijos nuokrypis neviršija specifikacijos ribos, o grotuvo garsas perduodamas per radijo siųstuvą. Architektūros peržiūros protokole nėra atvirų kritinių pastabų, įrangos sprendimas patvirtintas užsakymui. \\
    KT5 & Pagamintas ir patvirtintas pratimų vaizdo ir garso turinys & 2026-12-23 & Visų 40~pratimų ir 10~programų įrašai su garso instrukcijomis paruošti grotuvo formatu (Full HD); treneriai pasirašė turinio priėmimo lapą. \\
    KT6 & Veikia užsiėmimo vykdymo grandinė & 2027-01-11 & 1~iteracijos demonstracijoje užsiėmimo programa paleidžiama iš debesies per vykdymo variklio (S1) bandomąją sąsają, nes trenerio valdymo pultas (P2) kuriamas 3~iteracijoje. Grotuvas rodo pratimą ekrane, o trenerio balsas ir įrašytos instrukcijos girdimi ausinėse ir salės kolonėlėje; kritinių defektų nėra; pasirašytas 1~iteracijos demonstracijos protokolas. \\
    KT7 & Užsiėmimo programa vykdoma su filmuotu turiniu ir nutrūkus ryšiui & 2027-02-10 & Trenerio sudaryta užsiėmimo programa paleidžiama per S1 bandomąją sąsają ir įvykdoma ekrane ir ausinėse; nutrūkus ryšiui, grotuvas užbaigia programos etapą iš talpyklos. Pasirašytas 2~iteracijos demonstracijos protokolas. \\
    KT8 & Įranga sumontuota baseine ir sporto salėje & 2027-02-10 & Ekranai, grotuvai, radijo siųstuvas ir kolonėlės sumontuoti ir įregistruoti; montavimo aktas pasirašytas. \\
    KT9 & Veikia paskyros, užsiėmimų tvarkaraštis ir istorija & 2027-03-16 & Dalyvis prisiregistruoja, treneris suplanuoja užsiėmimą ir paleidžia jį iš valdymo pulto, istorija kaupiama. Pasirašytas 3~iteracijos demonstracijos protokolas. \\
    KT10 & Sistema funkciškai užbaigta ir patikrinta & 2027-04-16 & 4~iteracijos demonstracijoje veikia vietų rezervavimas, lankomumo žymėjimas ir asmeninių programų sudarymas; pasirašytas demonstracijos protokolas. Sistemos, vaidmenų prieigos ir saugumo bandymuose neliko kritinių ir aukšto prioriteto defektų. \\
    KT11 & Sistema veikia baseine ir sporto salėje & 2027-05-07 & Visi integracinių bandymų scenarijai atlikti realioje aplinkoje, protokolai pasirašyti, kritinių defektų nėra. \\
    KT12 & Treneriai ir administratoriai išmokyti & 2027-05-17 & Kiekvienas piloto treneris savarankiškai paleido ir užbaigė užsiėmimą; naudotojo vadovai perduoti, mokymų žiniaraštis pasirašytas. \\
    KT13 & Sistema priimta po pilotinio naudojimo & 2027-06-17 & Parengta 4~savaičių piloto ataskaita; užsakovas pasirašė naudotojų priėmimo bandymo protokolą be kritinių pastabų. \\
    KT14 & Sistema perduota užsakovui & 2027-06-29 & Pasirašytas priėmimo ir perdavimo aktas; perduoti išeities kodas, dokumentacija, prieigos duomenys ir įrangos inventorius. \\
    KT15 & Garantinis laikotarpis baigtas, projektas uždarytas & 2027-07-30 & Per 20~d.~d. garantinio laikotarpio kritiniai defektai ištaisyti; užsakovas patvirtino projekto uždarymo ataskaitą. \\
    \bottomrule
  \end{tabularx}
\end{table}

\subsection{Ganto diagrama}

Ganto diagrama (\ref{fig:gantt}~pav.) rodo projekto etapus, iteracijas, KT
ir lygiagrečius turinio bei įrangos darbus. Tamsios juostos žymi kritinio
kelio veiklas, šviesios – veiklas su laisvu laiku, punktyrinės – laiko
rezervus B1 ir B2, o rombai – KT. Diagrama sutraukta iki etapų ir iteracijų;
visos 75 veiklos su datomis ir priklausomybėmis pateiktos tvarkaraščio
faile \texttt{schedule.csv}. Įrangos grandinė (užsakymas, 30~d.~d.
pristatymas ir montavimas) turi 41~d.~d. laisvo laiko, todėl diagramoje
atskirta nuo kritinio kelio, o turinys paruošiamas dar prieš 2~iteraciją
(KT5 – anksčiau nei KT7).

\begin{landscape}
  \begin{figure}[htbp]
    \centering
    % Siūloma antraštė: Projekto tvarkaraštis. Tamsios juostos – kritinis kelias, šviesios – veiklos su laisvu laiku, punktyrinės – laiko rezervai, rombai – KT.
% Sugeneruota gantt.py iš datos.json; rankomis netaisyti.
\begin{ganttchart}[
  time slot format=isodate, x unit=0.058cm, y unit chart=0.31cm, y unit title=0.40cm,
  hgrid={*1{draw=black!8}}, title/.append style={draw=black!40, fill=black!4},
  title label font=\scriptsize, bar label font=\scriptsize, group label font=\scriptsize\bfseries,
  milestone label font=\scriptsize\itshape, bar height=0.55, group top shift=0.35, group height=0.3,
  group peaks height=0.15, group/.append style={fill=black!75}, milestone/.append style={fill=black},
  milestone height=0.7, milestone top shift=0.15, milestone left shift=-2, milestone right shift=3,
  link/.style={-latex, draw=black!70}
]{2026-10-01}{2027-08-03}
  \gantttitle{2026}{92} \gantttitle{2027}{215} \\
  \gantttitle{spal.}{31} \gantttitle{lapkr.}{30} \gantttitle{gruod.}{31} \gantttitle{saus.}{31} \gantttitle{vas.}{28} \gantttitle{kov.}{31} \gantttitle{bal.}{30} \gantttitle{geg.}{31} \gantttitle{birž.}{30} \gantttitle{liep.}{31} \gantttitle{}{3} \\
  \ganttgroup{Reikalavimų analizė}{2026-10-01}{2026-10-23} \\
  \ganttbar[bar/.append style={fill=black!65}]{Poreikių analizė ir reikalavimų specifikacija}{2026-10-01}{2026-10-23} \\
  \ganttmilestone{KT1 Pratimų ir programų sąrašas}{2026-10-15} \\
  \ganttmilestone{KT2 Reikalavimų specifikacija}{2026-10-23} \\
  \ganttgroup{Planavimas}{2026-10-26}{2026-11-04} \\
  \ganttbar[bar/.append style={fill=black!65}]{Projekto planas ir priėmimo kriterijai}{2026-10-26}{2026-11-04} \\
  \ganttmilestone{KT3 Projekto planas}{2026-11-04} \\
  \ganttgroup{Projektavimas}{2026-11-05}{2026-12-03} \\
  \ganttbar[bar/.append style={fill=black!65}]{Architektūra ir sąsajų prototipai}{2026-11-05}{2026-11-26} \\
  \ganttbar[bar/.append style={fill=black!65}]{Variklio, grotuvo ir garso prototipai}{2026-11-05}{2026-12-03} \\
  \ganttmilestone{KT4 Architektūra ir įrangos sprendimas}{2026-12-03} \\
  \ganttgroup{Kūrimas}{2026-12-04}{2027-04-09} \\
  \ganttbar[bar/.append style={fill=black!65}, name=it1]{1 iteracija}{2026-12-04}{2027-01-11} \\
  \ganttmilestone{KT6 Vykdymo grandinė}{2027-01-11} \\
  \ganttbar[bar/.append style={fill=black!65}, name=it2]{2 iteracija}{2027-01-08}{2027-02-10} \\
  \ganttmilestone{KT7 Programos su turiniu}{2027-02-10} \\
  \ganttbar[bar/.append style={fill=black!65}, name=it3]{3 iteracija}{2027-02-09}{2027-03-16} \\
  \ganttmilestone{KT9 Paskyros ir tvarkaraštis}{2027-03-16} \\
  \ganttbar[bar/.append style={fill=black!65}, name=it4]{4 iteracija}{2027-03-15}{2027-04-09} \\
  \ganttgroup{Diegimas ir palaikymas}{2027-04-12}{2027-07-30} \\
  \ganttbar[bar/.append style={fill=black!65}, name=sb]{Sistemos bandymai ir diegimas objekte}{2027-04-12}{2027-04-16} \\
  \ganttmilestone{KT10 Funkciškai užbaigta}{2027-04-16} \\
  \ganttbar[bar/.append style={fill=white, dashed}]{Laiko rezervas B1}{2027-04-19}{2027-04-28} \\
  \ganttbar[bar/.append style={fill=black!65}]{Integraciniai bandymai objekte}{2027-04-29}{2027-05-07} \\
  \ganttmilestone{KT11 Veikia baseine ir salėje}{2027-05-07} \\
  \ganttbar[bar/.append style={fill=black!65}]{Mokymai}{2027-05-10}{2027-05-17} \\
  \ganttmilestone{KT12 Mokymai baigti}{2027-05-17} \\
  \ganttbar[bar/.append style={fill=black!65}]{Pilotinis naudojimas}{2027-05-18}{2027-06-14} \\
  \ganttbar[bar/.append style={fill=black!65}]{Priėmimo bandymas}{2027-06-15}{2027-06-17} \\
  \ganttmilestone{KT13 Sistema priimta}{2027-06-17} \\
  \ganttbar[bar/.append style={fill=white, dashed}]{Laiko rezervas B2}{2027-06-18}{2027-06-25} \\
  \ganttbar[bar/.append style={fill=black!65}]{Perdavimas}{2027-06-28}{2027-06-29} \\
  \ganttmilestone{KT14 Sistema perduota}{2027-06-29} \\
  \ganttbar[bar/.append style={fill=black!65}]{Garantinis palaikymas (20 d. d.)}{2027-06-30}{2027-07-28} \\
  \ganttbar[bar/.append style={fill=black!65}]{Projekto uždarymas}{2027-07-29}{2027-07-30} \\
  \ganttmilestone{KT15 Projektas uždarytas}{2027-07-30} \\
  \ganttgroup[group/.append style={fill=black!30}]{Lygiagretūs darbai}{2026-10-08}{2027-02-10} \\
  \ganttbar[bar/.append style={fill=black!12}, name=turinys]{Turinys: filmavimas, garsas, montavimas}{2026-10-16}{2026-12-23} \\
  \ganttmilestone[name=kt5]{KT5 Turinys paruoštas}{2026-12-23} \\
  \ganttbar[bar/.append style={fill=black!12}]{Bandomasis įrangos komplektas}{2026-10-08}{2026-11-05} \\
  \ganttbar[bar/.append style={fill=black!12}]{Įrangos užsakymas ir pristatymas}{2026-12-03}{2027-01-21} \\
  \ganttbar[bar/.append style={fill=black!12}, name=montavimas]{Įrangos montavimas}{2027-01-22}{2027-02-10} \\
  \ganttmilestone{KT8 Įranga sumontuota}{2027-02-10}
  \ganttlink{kt5}{it2}
  \ganttlink{montavimas}{sb}
  \ganttlink{it4}{sb}
\end{ganttchart}

    \caption{Projekto tvarkaraštis. Tamsios juostos – kritinis kelias, šviesios – veiklos su laisvu laiku, punktyrinės – laiko rezervai, rombai – KT}
    \label{fig:gantt}
  \end{figure}
\end{landscape}

\subsection{Kritinis kelias ir laiko rezervai}

Kritinis kelias – ilgiausia vienas nuo kito priklausomų veiklų seka; ji
lemia trumpiausią įmanomą projekto trukmę, 208~d.~d. Kritiniame kelyje yra
reikalavimų analizė, projekto planas, architektūra ir prototipai su bandymu
baseine, visos keturios iteracijos, sistemos bandymai, bandymai objekte,
mokymai, pilotinis naudojimas, priėmimas, perdavimas ir garantinis
palaikymas. Šios veiklos neturi laisvo laiko, todėl bet kurios jų vėlavimas
nukelia projekto pabaigą.

Kitos veiklos turi laisvo laiko, t.~y. gali vėluoti nevėlindamos projekto
pabaigos (\ref{tab:laiko-rezervai}~lentelė). Pavyzdžiui, įrangos montavimas
gali baigtis iki 41~d.~d. vėliau (vėliausiai 2027~m. balandžio 13~d.), o
turinio gamyba – iki 11~d.~d. vėliau, ir perdavimo data nepasikeis.

Kritinį kelią saugo du laiko rezervai: B1 (8~d.~d. prieš bandymus objekte)
ir B2 (5~d.~d. prieš perdavimą). Per juos niekas nedirba, todėl sąnaudų jie
neturi. Jei kritinio kelio veikla vėluoja, pavyzdžiui, sistemos bandymai
užtrunka ilgiau arba baseinas laikinai neprieinamas, vėlavimą sugeria laiko
rezervas, ir svarbiausios datos nepasislenka. Laiko rezervų nereikia
painioti su sąnaudų rezervais (\ref{sec:dydis}~skyrius): pirmieji saugo
datas, antrieji – darbo dienų kiekį.

\begin{table}[H]
  \centering
  \caption{Laiko rezervai ir laisvas laikas}
  \label{tab:laiko-rezervai}
  \small
  \begin{tabularx}{\textwidth}{@{}>{\raggedright\arraybackslash}p{0.3\textwidth}>{\raggedright\arraybackslash}Xr@{}}
    \toprule
    \textbf{Pavadinimas} & \textbf{Paaiškinimas} & \textbf{D.~d.} \\
    \midrule
    Laiko rezervas~B1 & 2027-04-19~-- 2027-04-28, prieš integracinius bandymus objekte & 8 \\
    Laiko rezervas~B2 & 2027-06-18~-- 2027-06-25, prieš perdavimą & 5 \\
    \midrule
    Įrangos grandinės laisvas laikas & įrangos montavimas gali baigtis vėliausiai 2027-04-13 (planuojama 2027-02-10); tiekimo vėlavimas iki 41~d.~d. perdavimo nepaveiks & 41 \\
    Turinio grandinės laisvas laikas & turinys gali būti paruoštas vėliausiai 2027-01-12 (KT5 planuojamas 2026-12-23); KT7 data nepasikeis & 11 \\
    Naudotojo dokumentacijos laisvas laikas & naudotojo vadovai gali būti parengti vėliausiai 2027-05-07 (planuojama 2027-04-09), prieš mokymus & 20 \\
    \bottomrule
  \end{tabularx}
\end{table}

\subsection{Projekto stebėsena}

Projekto vadovas kas savaitę lygina planą su faktu. Jis tikrina, ar KT
pasiekti laiku ir ar įvykdyti jų kriterijai (\ref{tab:kt}~lentelė);
iteracijų pažanga matoma ir demonstracijose užsakovui. Kiekvienai veiklai
lyginamos planuotos ir faktinės sąnaudos (žm.~d.) ir vertinama, kiek darbo
iš tikrųjų atlikta, o ne tik kiek laiko praleista. Taip pat stebima, kiek
sunaudota laiko rezervų B1 ir B2 bei rizikų rezervo; rizikų rezervo
naudojimas susiejamas su konkrečia registro rizika pagal jos ID
(\ref{tab:rezervas}~lentelė).

Kas savaitę vyksta būsenos susirinkimas: aptariama, kas padaryta, kas bus
daroma toliau, kokios problemos trukdo ir kokių naujų rizikų atsirado.
Nustačius nukrypimą, įvertinama jo priežastis ir poveikis projektui; jei
reikia keisti planą, tai daroma pagal pakeitimų valdymo procedūrą.

\subsection{Planavimo prielaidos}

\begin{itemize}
  \item dirbama 5 dienas per savaitę po 8 valandas; valstybinės šventės yra
    nedarbo dienos (2026~m. lapkričio 2~d., gruodžio 24 ir 25~d., 2027~m.
    sausio 1~d., vasario 16~d., kovo 11~d., kovo 29~d., birželio 24~d. ir
    liepos 6~d.);
  \item veiklos trukmė apskaičiuojama sąnaudas padalijus iš joje dirbančių
    žmonių skaičiaus ir suapvalinus į didesnę pusę;
  \item įrangos pristatymas trunka 30~d.~d. nuo užsakymo; tai laukimo laikas,
    o ne sąnaudos;
  \item pratimai baseine filmuojami ne darbo valandomis, kai baseinas
    neužimtas;
  \item komandą sudaro 9 žmonės, kiekvienas atlieka vieną (sudėtinę) rolę;
    vienu metu dirba ne daugiau kaip 7.
\end{itemize}
